\section{Conclusion}
\label{sec:conclusion}

This paper reported the design, firmware implementation, an honest
object-oriented architecture audit, and a hardware-validated correctness
fix for HunStat2, our own low-cost AD5941/RP2040 potentiostat. The audit
found real but limited object orientation -- encapsulation in only two of
eight classes, no inheritance or polymorphism anywhere in the codebase,
and three method classes carrying verbatim-duplicated measurement-configuration
code -- and showed that this duplication was not a cosmetic concern: a
commercial electrochemical AFE removes most of the burden of getting a
low-cost potentiostat's analog performance right, but firmware can still
drive it into the wrong of its two excitation/sensing loops, or
misinterpret the numeric convention of the code it returns, and produce
output that compiles cleanly and runs without crashing while being
physically wrong -- and when that firmware is triplicated, so is the bug.
We diagnosed two such defects in HunStat2's chronoamperometry,
square-wave voltammetry, and differential pulse voltammetry firmware -- an
incorrect High-Speed-loop configuration where every Analog Devices
DC-method reference example uses the Low-Power loop, and an incorrect ADC
zero-point convention -- by direct, file-by-file comparison against the
vendor's own reference examples, the most reliable diagnostic method
available when the correct behavior is already fully specified in code the
manufacturer publishes. Both defects, together with an undersized ADC
polling-timeout budget sharing the same root cause, were corrected
identically across all three affected classes, once per copy. The fix is
carried forward into the same main firmware independently validated by
two companion manuscripts: their data confirms CA, CV, SWV,
and DPV all produce correctly-signed, signal-dependent, non-degenerate
output, and the companion metrology manuscript's traceable accuracy check
($R^2\geq0.999994$, $+0.428\,\%$ recovered-resistance error) is
independent confirmation that the corrected signal chain is quantitatively
correct, not merely internally consistent. We also mapped the remaining
differences from the vendor reference into deliberate architectural
choices appropriate to this instrument's live, host-driven use case, and
genuine, unreconciled gaps -- RTIA live-calibration and sample-timing
precision chief among them. This paper deliberately stops short of
restructuring the duplication its own audit documents; that refactor, and
its own hardware-measured proof that it leaves every method's output
unchanged, is the companion architecture manuscript's
contribution~\cite{clement2026architecture}. Read together with that
manuscript and the companion metrology manuscript's traceability
protocol~\cite{clement2026linearity}, this instrument family now has an
accurate map of what is fixed, what is a considered trade-off, and what is
still open.
